\section{Shell Profiles at Split Finite Places}
\label{fw:section}

We now place locally constant weights at finitely many finite places of
$F$ that split in $K$; in Section~\ref{sec:certificate} these are the
places above the selected primes of $B$ (Definition~\ref{an:selected-def}).
In Lemma~\ref{geo:selection}, the count at each prime of $F$ above
$r\in\mathcal R$ corresponds to the indicator function of a product of
two balls in the completion, the unweighted shell profile of
Definition~\ref{fw:shell-profile}; shell profiles replace this indicator
by weighted sums of indicators of products of shells
(Corollary~\ref{fw:uniform-types}). The main result is
Proposition~\ref{fw:transfer}, the transfer with shell profiles, which is the
form of the construction applied in Section~\ref{sec:certificate}; its
margin \eqref{fw:margin} replaces the terms $J-\delta H$ of
\eqref{geo:general-margin} by the values of a local functional
(Definition~\ref{fw:local-functional}). We keep the hypotheses (G1)--(G3)
and the notation of Section~\ref{geo:section}, and put
\[
 \covol_0=\covol\bigl((\mathcal O_K)_\infty\bigr)=(\lambda/2)^d .
\]

\subsection{The Local Functional}\label{fw:local-subsection}

This subsection defines the local functional at a split finite place and
computes it for shell profiles (Lemma~\ref{fw:shell-lemma}); its values are
the terms $\log\mathcal F_{\delta,\mathfrak p}$ of the margin
\eqref{fw:margin}.

Let $\mathsf L$ be a nonarchimedean local field with valuation ring $\mathcal
O$, uniformizer $\varpi$ and residue field of cardinality $Q$. We use
the additive Haar measure on $\mathsf L$ with $|\mathcal O|=1$, the product
measure on the plane $\mathsf L^2=\mathsf L\times\mathsf L$, and the
multiplicative Haar measure $d\omega$ on $\mathcal O^\times$ of total
measure one.

\begin{definition}\label{fw:shift-operator}
For $n\in\Z$ and $\omega\in\mathcal O^\times$ put
\[
 s(n,\omega)=(\varpi^n\omega,\ \varpi^{-n}\omega^{-1})\in\mathsf L^2,\qquad
 \mathcal Tg(z)=\sum_{n\in\Z}\int_{\mathcal O^\times}g(z+s(n,\omega))\,d\omega
\]
for a locally constant, compactly supported function $g$ on $\mathsf L^2$.
\end{definition}

The two coordinates of $s(n,\omega)$ have product one; they model the
finite coordinates of an element of relative norm one at a split place.

\begin{lemma}\label{fw:finite-sum}
Let $\varphi_1$ and $\varphi_2$ be locally constant and compactly supported
on $\mathsf L^2$. Then only finitely many $n$ contribute to
$\int_{\mathsf L^2}\varphi_1\,\mathcal T\varphi_2$, which is therefore a
finite sum of finite integrals.
\end{lemma}

\begin{proof}
If $\varphi_1(z)\varphi_2(z+s(n,\omega))\ne0$, then both coordinates of
$s(n,\omega)$ lie in the compact set
$\operatorname{supp}\varphi_2-\operatorname{supp}\varphi_1$, which bounds
$n$ from above and from below.
\end{proof}

\begin{definition}\label{fw:local-functional}
For locally constant, compactly supported functions $\varphi_1,\varphi_2$ on
$\mathsf L^2$ put $\langle\varphi_1,\varphi_2\rangle_{\mathcal
T}=\int_{\mathsf L^2}\varphi_1\,\mathcal T\varphi_2$
(Lemma~\ref{fw:finite-sum}). For a nonnegative, locally constant, compactly
supported function $g$ on $\mathsf L^2$ with $g\ne0$, and $0<\delta<1$,
$p=2/(1+\delta)$, put
\begin{equation}\label{fw:functional}
 A(g)=\int_{\mathsf L^2}g^p,\qquad I(g)=\langle g,g\rangle_{\mathcal T},\qquad
 \mathcal F_\delta(g)=\frac{I(g)}{A(g)^{1+\delta}} .
\end{equation}
We call $A(g)$ the \emph{mass}, $I(g)$ the \emph{overlap} and
$\mathcal F_\delta$ the \emph{local functional}, as for the archimedean
masses and overlaps of Definition~\fref{geo:profiles}.
\end{definition}

For $i\in\Z$ let $\mathcal B_i=\varpi^{-i}\mathcal O$, a ball of measure
$Q^i$. Write $[x]_+=\max(x,0)$.

\begin{lemma}\label{fw:gram-lemma}
For all integers $i_1,j_1,i_2,j_2$,
\begin{equation}\label{fw:rectangle-gram}
 \bigl\langle\mathbf 1_{\mathcal B_{i_1}\times\mathcal B_{j_1}},
 \mathbf 1_{\mathcal B_{i_2}\times\mathcal B_{j_2}}\bigr\rangle_{\mathcal T}
 =Q^{\min(i_1,i_2)+\min(j_1,j_2)}\,
 [\max(i_1,i_2)+\max(j_1,j_2)+1]_+.
\end{equation}
\end{lemma}

\begin{proof}
Two balls of an ultrametric space are disjoint or nested. Hence
$\mathcal B_{i_1}\cap(\mathcal B_{i_2}-x)$ has measure
$Q^{\min(i_1,i_2)}$ if $x\in\mathcal B_{\max(i_1,i_2)}$, and is empty
otherwise. The left side of \eqref{fw:rectangle-gram} is therefore
$Q^{\min(i_1,i_2)+\min(j_1,j_2)}$ times the measure of the set of
$(n,\omega)$ with $\varpi^n\omega\in\mathcal B_{\max(i_1,i_2)}$ and
$\varpi^{-n}\omega^{-1}\in\mathcal B_{\max(j_1,j_2)}$. These conditions
say $-\max(i_1,i_2)\le n\le\max(j_1,j_2)$ and do not involve $\omega$,
and $d\omega$ has total measure one.
\end{proof}

\begin{definition}\label{fw:shell-profile}
Put $\mathcal S_0=\mathcal O$ and $\mathcal S_i=\mathcal B_i\setminus
\mathcal B_{i-1}$ for $i\ge1$. These are the \emph{shells}, and their
measures $m_0=1$ and $m_i=Q^i-Q^{i-1}$ are the \emph{shell measures}.
Let $k\ge0$ be an integer, and let $w=(w_{ij})_{i,j\ge0}$ be nonnegative
weights, finitely many of them nonzero, with $w_{00}>0$. The
\emph{shell profile} with parameters $(Q,k,w)$ on $\mathsf L^2$ is
\[
 g=\sum_{i,j\ge0}w_{ij}\mathbf 1_{\mathcal S_i\times\varpi^{-k}\mathcal S_j}.
\]
The shell profile with $w_{00}=1$ and $w_{ij}=0$ otherwise, that is, the
indicator function of $\mathcal O\times\varpi^{-k}\mathcal O$, is the
\emph{unweighted shell profile} with parameters $(Q,k)$.
\end{definition}

\begin{lemma}\label{fw:shell-invariance}
A shell profile $g$ with parameters $(Q,k,w)$ is nonnegative, locally
constant and compactly supported. It is invariant under
$(x,y)\mapsto(\omega x,\omega'y)$ for all $\omega,\omega'\in\mathcal
O^\times$; in particular it is even. It is also invariant under
translation by the group
\begin{equation}\label{fw:hard-period}
 \mathcal C=\mathcal O\times\varpi^{-k}\mathcal O,\qquad |\mathcal C|=Q^k.
\end{equation}
\end{lemma}

\begin{proof}
The first assertion holds because the shells are compact open sets and
finitely many of the nonnegative weights are nonzero. Multiplication by
a unit preserves valuations and hence each shell. Adding an element of
$\mathcal O$ does not change the shell of a point, whatever the number
of shells; this gives the invariance under $\mathcal C$, applied to the
first coordinate and to $\varpi^k$ times the second.
\end{proof}

We call $\mathcal C$ the \emph{period} of $g$.

\begin{definition}\label{fw:gram-matrix}
Let $Q$ and $k$ be as in Definition~\ref{fw:shell-profile}, with shell
measures $m_i$. The symmetric matrix $\mathbf D=(\mathbf D_{ii'})_{i,i'\ge0}$
is defined by
\[
 \mathbf D_{ii'}=m_{\min(i,i')}\ (i\ne i'),\qquad \mathbf D_{00}=0,\qquad
 \mathbf D_{ii}=im_i-Q^{i-1}\ (i\ge1),
\]
and the \emph{shell Gram matrix} $\mathbf G$, indexed by pairs of pairs of
nonnegative integers, by
\[
 \mathbf G_{(i,j),(i',j')}=(k+1)m_im_j\mathbf 1_{i=i',\,j=j'}
 +m_j\mathbf D_{ii'}\mathbf 1_{j=j'}+m_i\mathbf D_{jj'}\mathbf 1_{i=i'}.
\]
\end{definition}

\begin{lemma}\label{fw:shell-lemma}
The shell profile $g$ with parameters $(Q,k,w)$ satisfies, with the shell
measures $m_i$ and the shell Gram matrix $\mathbf G$ of
Definition~\ref{fw:gram-matrix},
\begin{equation}\label{fw:shell-kernel}
 A(g)=Q^k\sum_{i,j}m_im_jw_{ij}^p,\qquad
 I(g)=Q^k\sum_{(i,j),(i',j')}w_{ij}\,\mathbf G_{(i,j),(i',j')}\,w_{i'j'} .
\end{equation}
All entries of $\mathbf G$ are nonnegative, and
$I(g)\ge(k+1)Q^kw_{00}^2>0$. In particular $\mathcal F_\delta(g)$
depends only on the parameters, and the unweighted shell profile has
$\mathcal F_\delta(g)=(k+1)Q^{-\delta k}$.
\end{lemma}

\begin{proof}
The formula for $A(g)$ holds because the sets
$\mathcal S_i\times\varpi^{-k}\mathcal S_j$ are disjoint of measure
$Q^km_im_j$. For $I(g)$, write $\mathbf 1_{\mathcal S_0}=\mathbf
1_{\mathcal B_0}$ and $\mathbf 1_{\mathcal S_i}=\mathbf 1_{\mathcal
B_i}-\mathbf 1_{\mathcal B_{i-1}}$ for $i\ge1$, and note
$\varpi^{-k}\mathcal B_j=\mathcal B_{j+k}$. Only balls with nonnegative
indices occur, so by \eqref{fw:rectangle-gram} and $k\ge0$ the positive
part can be dropped, and
$\langle\mathbf 1_{\mathcal B_{i_1}\times\mathcal B_{j_1+k}},
\mathbf 1_{\mathcal B_{i_2}\times\mathcal B_{j_2+k}}\rangle_{\mathcal T}$,
for $i_1,j_1,i_2,j_2\ge0$, is
\[
 Q^k\bigl[(k+1)Q^{\min(i_1,i_2)}Q^{\min(j_1,j_2)}
 +\max(i_1,i_2)Q^{\min(i_1,i_2)}Q^{\min(j_1,j_2)}
 +Q^{\min(i_1,i_2)}\max(j_1,j_2)Q^{\min(j_1,j_2)}\bigr].
\]
For a function $\varphi$ of two nonnegative integers put
$(\partial\varphi)(i,i')=\varphi(i,i')-\varphi(i-1,i')-\varphi(i,i'-1)
+\varphi(i-1,i'-1)$, where terms with an index $-1$ are omitted; this is
the expansion of the shells into balls. Apply $\partial$ in
$(i_1,i_2)$ and in $(j_1,j_2)$. A direct computation gives, for
$i,i'\ge0$, that $\partial Q^{\min(i_1,i_2)}$ at $(i,i')$ is
$m_i\mathbf 1_{i=i'}$, and that
$\partial\bigl(\max(i_1,i_2)Q^{\min(i_1,i_2)}\bigr)$ is $\mathbf D_{ii'}$.
For example, when $1\le i<i'$ the latter is
$i'Q^i-i'Q^{i-1}-(i'-1)Q^i+(i'-1)Q^{i-1}=m_i$, and when $i=i'\ge1$ it is
$iQ^i-2iQ^{i-1}+(i-1)Q^{i-1}=im_i-Q^{i-1}$. This gives
$\langle\mathbf 1_{\mathcal S_i\times\varpi^{-k}\mathcal S_j},
\mathbf 1_{\mathcal S_{i'}\times\varpi^{-k}\mathcal S_{j'}}\rangle_{\mathcal T}
=Q^k\mathbf G_{(i,j),(i',j')}$, and hence the formula for $I(g)$.
All entries of $\mathbf G$ are nonnegative, since
$\mathbf D_{ii}=Q^{i-1}(iQ-i-1)\ge0$ for $i\ge1$; as
$\mathbf G_{(0,0),(0,0)}=k+1$, this gives $I(g)\ge(k+1)Q^kw_{00}^2$.
For the unweighted shell profile only $\mathbf G_{(0,0),(0,0)}=k+1$
contributes.
\end{proof}

\begin{definition}\label{fw:endpoint-law}
Let $g$ be a shell profile. By Lemma~\ref{fw:shell-lemma}, $I(g)>0$, and by
the definition of $I(g)$ the function
$(z,n,\omega)\mapsto g(z)g(z+s(n,\omega))/I(g)$ is a probability density on
$\mathsf L^2\times\Z\times\mathcal O^\times$, with counting measure on $\Z$.
As at the archimedean places (Definition~\fref{geo:profiles}), we call it the
\emph{endpoint law} of $g$; its first endpoint is $z$ and its second endpoint
is $z+s(n,\omega)$.
\end{definition}

\begin{corollary}\label{fw:product-weights}
Let $x=(x_i)_{i\ge0}$ and $x'=(x'_j)_{j\ge0}$ be finitely supported
vectors of nonnegative numbers with $x_0x'_0>0$, and with $\mathbf D$ as in
Definition~\ref{fw:gram-matrix} put
\[
 P_x=\sum_im_ix_i^p,\qquad S_x=\sum_im_ix_i^2,\qquad
 R_x=\sum_{i,i'}x_i\mathbf D_{ii'}x_{i'},
\]
and define $P_{x'},S_{x'},R_{x'}$ in the same way. The shell profile
$g$ with parameters $(Q,k,w)$ and product weights $w_{ij}=x_ix'_j$
satisfies
\begin{equation}\label{fw:product-shells}
 A(g)=Q^kP_xP_{x'},\qquad
 I(g)=Q^k\bigl[(k+1)S_xS_{x'}+R_xS_{x'}+S_xR_{x'}\bigr].
\end{equation}
\end{corollary}

\begin{proof}
Substitute $w_{ij}=x_ix'_j$ in \eqref{fw:shell-kernel}.
\end{proof}

\begin{remark}\label{fw:one-shell}
Suppose that only $\mathcal S_0$ and $\mathcal S_1$ carry weight, with
$x=x'=(1,t)$ and $t\ge0$. Then $S_x=1+(Q-1)t^2$ and
$R_x=2t+(Q-2)t^2$, and the difference between $\log\mathcal F_\delta(g)$
and the value $\log\bigl((k+1)Q^{-\delta k}\bigr)$ of the unweighted
shell profile is
\begin{equation}\label{fw:one-shell-gain}
 \log\frac{(k+1)S_x^2+2S_xR_x}{k+1}-2(1+\delta)\log\bigl(1+(Q-1)t^p\bigr).
\end{equation}
Since $p>1$, its right derivative at $t=0$ is $4/(k+1)>0$, so a
sufficiently small weight on $\mathcal S_1$ always improves on the
unweighted shell profile.
\end{remark}

\begin{remark}\label{fw:general-weights}
Section~\ref{sec:certificate} uses general weights $w_{ij}$, symmetric
$10\times10$ matrices, at the places above $2$, $3$ and $5$, through the
formulas \eqref{fw:shell-kernel}, and product weights, through
\eqref{fw:product-shells}, at the places above $29$ and $\mathfrak t_3$.
Nothing in the rest of this section, in particular in the proof of
Proposition~\ref{fw:transfer}, uses the product form: it uses only that a
shell profile is nonnegative, locally constant and compactly supported,
invariant under units in each coordinate and under its period
(Lemma~\ref{fw:shell-invariance}), with $I(g)>0$
(Lemma~\ref{fw:shell-lemma}), and that its parameters range over a finite
set. As a check of \eqref{fw:shell-kernel} for general weights, the
supplementary program \texttt{general\_windows\_check.py} compares it with a
direct computation in a finite model of the local field, in $48$ cases, $32$ of
them with random weights and $16$ with a fixed family of monotone weights that
are not symmetric, and finds exact agreement; the program \texttt{general\_windows.py} that evaluates
\eqref{fw:shell-kernel} for the certificate reproduces
\eqref{fw:product-shells} exactly for product weights.
\end{remark}

\subsection{Finite Periods and Lattice Counting}

This subsection adds finite coordinates at the places of a finite set
$\mathcal P$ of split places and proves the lattice-point bound of
Lemma~\ref{fw:counting}, the analogue of Lemma~\ref{geo:poisson}.

Let $\mathcal P$ be a finite set of finite places of $F$ that split in
$K$. For $\mathfrak p\in\mathcal P$ write $\mathfrak p\mathcal
O_K=\mathfrak P_{\mathfrak p}\,\iota\mathfrak P_{\mathfrak p}$, where we
also write $\mathfrak p$ for the prime ideal of $F$. The completion of
$K$ at $\mathfrak P_{\mathfrak p}$ is $F_{\mathfrak p}$; for $x\in K$ let
$x_{\mathfrak p}\in F_{\mathfrak p}$ be its image under the completion
map. The map $x\mapsto(x_{\mathfrak p},(\iota x)_{\mathfrak p})$ records
the completions at $\mathfrak P_{\mathfrak p}$ and at $\iota\mathfrak
P_{\mathfrak p}$, and $\iota$ exchanges the two coordinates. If
$N_{K/F}\beta=1$, then $\beta_{\mathfrak p}(\iota\beta)_{\mathfrak p}=1$.
Let $\mathcal O_{\mathfrak p}$, $\varpi_{\mathfrak p}$ and $Q_{\mathfrak
p}$ be the valuation ring, a uniformizer and the residue cardinality of
$F_{\mathfrak p}$; thus $Q_{\mathfrak p}=\mathrm N\mathfrak p$, the absolute norm of
$\mathfrak p$, and $Q_{\mathfrak p}$ plays the role of $Q$ in
Subsection~\ref{fw:local-subsection} for $\mathsf L=F_{\mathfrak p}$.

\begin{definition}\label{fw:finite-coordinates}
Put
\[
 \mathbb A_{\mathcal P}=\prod_{\mathfrak p\in\mathcal P}
 (F_{\mathfrak p}\times F_{\mathfrak p}),
\]
with the Haar measure that gives $\prod_{\mathfrak p}(\mathcal
O_{\mathfrak p}\times\mathcal O_{\mathfrak p})$ measure one. For a point
$(z,\zeta)$ of $\C^d\times\mathbb A_{\mathcal P}$ we call $z$ its
\emph{archimedean part} and $\zeta$ its \emph{finite part}. Let
$\mathcal O_{\mathcal P}=\mathcal O_{K,\mathcal P}$ be the ring of elements of $K$
that are integral at every prime other than the $\mathfrak P_{\mathfrak p}$
and $\iota\mathfrak P_{\mathfrak p}$. We embed $\mathcal O_{\mathcal P}$ in
$\C^d\times\mathbb A_{\mathcal P}$ by
$\gamma\mapsto(\gamma_\infty,\gamma_f)$, where
$\gamma_f=(\gamma_{\mathfrak p},(\iota\gamma)_{\mathfrak p})_{\mathfrak
p}$; thus $\gamma_\infty$ and $\gamma_f$ are the archimedean part and the
finite part of the image of $\gamma$. A subgroup of $\mathbb A_{\mathcal
P}$ of the form
\[
 \mathcal C_f=\prod_{\mathfrak p\in\mathcal P}\bigl(\varpi_{\mathfrak
 p}^{a_{\mathfrak p}}\mathcal O_{\mathfrak p}\times\varpi_{\mathfrak
 p}^{b_{\mathfrak p}}\mathcal O_{\mathfrak p}\bigr),\qquad
 a_{\mathfrak p},b_{\mathfrak p}\in\Z,
\]
is called \emph{rectangular}.
\end{definition}

A rectangular group has measure
$|\mathcal C_f|=\prod_{\mathfrak p}Q_{\mathfrak p}^{-a_{\mathfrak
p}-b_{\mathfrak p}}$.

\begin{lemma}\label{fw:lattice}
Let $\mathcal C_f$ be rectangular and put
$\mathfrak a_{\mathcal C_f}=\prod_{\mathfrak p}\mathfrak P_{\mathfrak p}^{a_{\mathfrak
p}}(\iota\mathfrak P_{\mathfrak p})^{b_{\mathfrak p}}$.
\begin{enumerate}
\item[(a)] The elements $\gamma\in\mathcal O_{\mathcal P}$ with $\gamma_f\in\mathcal C_f$
form $\mathfrak a_{\mathcal C_f}$, and $\mathrm N\mathfrak a_{\mathcal C_f}=|\mathcal C_f|^{-1}$, so
$\covol((\mathfrak a_{\mathcal C_f})_\infty)=\covol_0/|\mathcal C_f|$.
\item[(b)] Every coset of $\mathcal C_f$ in $\mathbb A_{\mathcal P}$
contains $\gamma_f$ for some $\gamma\in\mathcal O_{\mathcal P}$.
\item[(c)] $\mathcal O_{\mathcal P}$ is a lattice in $\C^d\times\mathbb A_{\mathcal P}$,
that is, a discrete subgroup with a fundamental domain of finite
measure. If $P_{\mathcal C_f}$ is a fundamental parallelepiped of
$(\mathfrak a_{\mathcal C_f})_\infty$, then $P_{\mathcal C_f}\times\mathcal C_f$ is a
fundamental domain for $\mathcal O_{\mathcal P}$, of measure $\covol_0$.
\item[(d)] For $h\in\R^c$, the set
$\mathcal O_{\mathcal P,h}=\{(W_h\gamma_\infty,\gamma_f):\gamma\in\mathcal
O_{\mathcal P}\}$ is again a lattice, with fundamental domain $W_hP_{\mathcal C_f}\times\mathcal C_f$
of measure $\covol_0$. Put $H_f=d^{-1}\log|\mathcal C_f|$. Every nonzero
vector $\xi$ of the dual lattice of $W_h(\mathfrak a_{\mathcal C_f})_\infty$ satisfies
$\prod_w|\xi_w|\ge\mu_f^d$, where $\mu_f=2e^{H_f/2-\ell}$.
\end{enumerate}
\end{lemma}

\begin{proof}
(a) The valuation of $\gamma_{\mathfrak p}$ is the exponent of
$\mathfrak P_{\mathfrak p}$ in $\gamma$, and that of $(\iota\gamma)_{\mathfrak
p}$ is the exponent of $\iota\mathfrak P_{\mathfrak p}$. Together with
integrality at the other primes, $\gamma_f\in\mathcal C_f$ says exactly
that $\gamma\in\mathfrak a_{\mathcal C_f}$. As $\mathrm N\mathfrak P_{\mathfrak
p}=\mathrm N(\iota\mathfrak P_{\mathfrak p})=Q_{\mathfrak p}$, we get
$\mathrm N\mathfrak a_{\mathcal C_f}=\prod_{\mathfrak p}Q_{\mathfrak p}^{a_{\mathfrak
p}+b_{\mathfrak p}}=|\mathcal C_f|^{-1}$, and the covolume follows from
the covolume formula \eqref{geo:covolume}.

(b) Choose $\tilde a_{\mathfrak p}\le a_{\mathfrak p}$ and $\tilde
b_{\mathfrak p}\le b_{\mathfrak p}$ such that the given coset lies in
$\tilde{\mathcal C}_f=\prod_{\mathfrak p}(\varpi_{\mathfrak p}^{\tilde
a_{\mathfrak p}}\mathcal O_{\mathfrak p}\times\varpi_{\mathfrak
p}^{\tilde b_{\mathfrak p}}\mathcal O_{\mathfrak p})$, and let
$\mathfrak a_{\tilde{\mathcal C}_f}\supset\mathfrak a_{\mathcal C_f}$ be the corresponding ideal. The
map $\gamma\mapsto\gamma_f$ induces a homomorphism $\mathfrak a_{\tilde{\mathcal C}_f}/\mathfrak a_{\mathcal C_f}\to\tilde{\mathcal C}_f/\mathcal C_f$. It is injective by
(a), applied to $\mathcal C_f$. Both groups are finite of the same
order: $|\tilde{\mathcal C}_f/\mathcal C_f|=|\tilde{\mathcal
C}_f|/|\mathcal C_f|$, and $|\mathfrak a_{\tilde{\mathcal C}_f}/\mathfrak a_{\mathcal C_f}|=\mathrm N\mathfrak a_{\mathcal C_f}/\mathrm N\mathfrak a_{\tilde{\mathcal C}_f}$ \cite[Proposition~4.2]{MilneANT}, which is the
same number by (a). So the map is bijective, and every coset of
$\mathcal C_f$ in $\tilde{\mathcal C}_f$ contains $\gamma_f$ for some
$\gamma\in\mathfrak a_{\tilde{\mathcal C}_f}\subset\mathcal O_{\mathcal P}$. No principality is needed.

(c) Given $(z,\zeta)\in\C^d\times\mathbb A_{\mathcal P}$, by (b) there is
$\gamma\in\mathcal O_{\mathcal P}$ with $\zeta-\gamma_f\in\mathcal C_f$, and then a unique
$\gamma'\in\mathfrak a_{\mathcal C_f}$ with $z-(\gamma+\gamma')_\infty\in P_{\mathcal C_f}$. If $(z,\zeta)$ had two such representations, the difference of the
two elements of $\mathcal O_{\mathcal P}$ would have finite part in $\mathcal C_f$, hence
lie in $\mathfrak a_{\mathcal C_f}$ by (a), and then it would vanish because
$P_{\mathcal C_f}$ is a fundamental domain. Thus
$P_{\mathcal C_f}\times\mathcal C_f$ is a fundamental domain, of measure
$(\covol_0/|\mathcal C_f|)\cdot|\mathcal C_f|=\covol_0$. Finally $\mathcal O_{\mathcal P}$
is discrete: for a bounded open set $Y\ni0$ in $\C^d$, its intersection
with the neighborhood $Y\times\mathcal C_f$ of $0$ consists of the
finitely many $\gamma\in\mathfrak a_{\mathcal C_f}$ with $\gamma_\infty\in Y$.

(d) The map $(z,\zeta)\mapsto(W_hz,\zeta)$ preserves measure, because
$W_h$ is real diagonal of determinant one, and it carries $\mathcal O_{\mathcal P}$ to
$\mathcal O_{\mathcal P,h}$ and $P_{\mathcal C_f}\times\mathcal C_f$ to $W_hP_{\mathcal C_f}\times\mathcal C_f$; so the first assertion follows from (c). By (a),
$\mathrm N\mathfrak a_{\mathcal C_f}=e^{-dH_f}$, and Lemma~\ref{geo:dual} gives the bound for
the dual vectors.
\end{proof}

Accordingly, the Fourier condition used in this section is
\begin{equation}\label{fw:fourier-condition}
 \varsigma\cdot2e^{H_f/2-\ell}\ge\log M+2\log5+2 .
\end{equation}

\begin{lemma}\label{fw:counting}
Let $(f_\R,g_\C)$ be admissible with Fourier constants $(M,\varsigma)$
(Definition~\ref{geo:admissible}), and
let $\Psi_\infty=\prod_vf_\R^p\prod_jg_\C^p$ on $\C^d$, with one factor
for each real place $v$ and each complex place $j$ of $F$
(Lemma~\ref{geo:tensor}). Let $\Psi_f\ge0$ be a function on $\mathbb
A_{\mathcal P}$ that is invariant under translation by a rectangular
group $\mathcal C_f$ and vanishes outside a finite union of its cosets.
If \eqref{fw:fourier-condition} holds, with $H_f$ defined by this
$\mathcal C_f$, then for every $h\in\R^c$ and every
$z_0\in\C^d\times\mathbb A_{\mathcal P}$ (a point of the whole space,
not an archimedean part),
\begin{equation}\label{fw:weighted-count}
 \sum_{(z,\zeta)\in\mathcal O_{\mathcal P,h}+z_0}\Psi_\infty(z)\Psi_f(\zeta)
 \le\frac2{\covol_0}\int\Psi_\infty\int\Psi_f .
\end{equation}
\end{lemma}

\begin{proof}
Group the points by the coset of $\mathcal C_f$ containing their finite
part. If two points of $\mathcal O_{\mathcal P,h}+z_0$ have finite parts in the same
coset, their difference comes from an element of $\mathfrak a_{\mathcal C_f}$, by
Lemma~\ref{fw:lattice}(a). So the archimedean parts of the points in one
coset form a translate of $W_h(\mathfrak a_{\mathcal C_f})_\infty$, or the empty set. By
Lemma~\ref{geo:poisson}, with the dual bound of
Lemma~\ref{fw:lattice}(d) and the Fourier bound for $\Psi_\infty$ of
Lemma~\ref{geo:tensor}, the sum of $\Psi_\infty$ over such a translate
is at most $2|\mathcal C_f|\int\Psi_\infty/\covol_0$. The values of
$\Psi_f$ on its finitely many nonzero cosets sum to
$\int\Psi_f/|\mathcal C_f|$.
\end{proof}

\begin{lemma}\label{fw:rescaling}
Let $\alpha\in\mathbb A_{\mathcal P}$ have coordinates
$(\varpi_{\mathfrak p}^{n_{\mathfrak p}},\varpi_{\mathfrak
p}^{-n_{\mathfrak p}})$ with $n_{\mathfrak p}\in\Z$. If $\Psi_f$ satisfies the
hypotheses of Lemma~\ref{fw:counting} with a rectangular group $\mathcal
C_f$, then $\Psi_f\circ\alpha$ satisfies them with the rectangular group
$\alpha^{-1}\mathcal C_f$, which gives the same $H_f$, and
$\int\Psi_f\circ\alpha=\int\Psi_f$.
\end{lemma}

\begin{proof}
Multiplication by $\alpha$ preserves the Haar measure, since at each
$\mathfrak p$ it scales the two coordinates by factors of absolute values
$Q_{\mathfrak p}^{-n_{\mathfrak p}}$ and $Q_{\mathfrak p}^{n_{\mathfrak p}}$.
For a rectangular group $\mathcal C_f$ the group $\alpha^{-1}\mathcal C_f$ is
rectangular, with exponents $a_{\mathfrak p}-n_{\mathfrak p}$ and
$b_{\mathfrak p}+n_{\mathfrak p}$, and of the same measure. The function
$\Psi_f\circ\alpha$ is invariant under translation by $\alpha^{-1}\mathcal
C_f$ and vanishes outside the images under $\alpha^{-1}$ of finitely many
cosets of $\mathcal C_f$, which are cosets of $\alpha^{-1}\mathcal C_f$.
\end{proof}

\subsection{Class Selection}

This subsection is the analogue of Lemma~\ref{geo:selection} for weighted
counts: it selects one ideal class and produces elements of relative norm one
with prescribed valuations at the places of $\mathcal P$
(Lemma~\ref{fw:selection}).

For $u\in\R^c$ and $n=(n_{\mathfrak p})\in\Z^{\mathcal P}$ let
$\mathbf e(u,n)\in\C^d\times\mathbb A_{\mathcal P}$ have archimedean part
the unit-norm displacement $\mathbf e(u)$ of Definition~\ref{geo:window} and
finite part $(\varpi_{\mathfrak p}^{n_{\mathfrak
p}},\varpi_{\mathfrak p}^{-n_{\mathfrak p}})_{\mathfrak p}$.

\begin{lemma}\label{fw:selection}
Let $\mathcal N\subset\Z^{\mathcal P}$ be finite, and let
$\mathcal W\colon\mathcal N\to[0,\infty)$ have positive sum. There are a
subset $\mathcal N'\subset\mathcal N$ with
\[
 \sum_{n\in\mathcal N'}\mathcal W(n)\ge\frac1{h_{\rm cap}h_{\rm rel}}
 \sum_{n\in\mathcal N}\mathcal W(n),
\]
an element $n_0\in\mathcal N'$ with $\mathcal W(n_0)>0$, and elements
$\beta_n\in\mathcal O_{\mathcal P}$, for $n\in\mathcal N'$, with $N_{K/F}\beta_n=1$ and
\[
 \beta_n\mathcal O_K=\prod_{\mathfrak p}\mathfrak P_{\mathfrak
 p}^{\,n_{\mathfrak p}-n_{0,\mathfrak p}}
 (\iota\mathfrak P_{\mathfrak p})^{-(n_{\mathfrak p}-n_{0,\mathfrak p})}.
\]
The cosets $\beta_n\mathcal O_K^1$ are pairwise disjoint, and
$\beta_n\mathcal O_K^1$ is the set of all elements of relative norm one
that generate this ideal. Moreover, let $\varpi_0\in\mathbb A_{\mathcal
P}$ have coordinates $(\varpi_{\mathfrak p}^{n_{0,\mathfrak
p}},\varpi_{\mathfrak p}^{-n_{0,\mathfrak p}})$. Then for $n\in\mathcal
N'$ and $\beta\in\beta_n\mathcal O_K^1$, the point $\varpi_0\beta_f$ has
coordinates $(\varpi_{\mathfrak p}^{n_{\mathfrak p}}\omega_{\mathfrak
p},\varpi_{\mathfrak p}^{-n_{\mathfrak p}}\omega_{\mathfrak p}^{-1})$
with units $\omega_{\mathfrak p}\in\mathcal O_{\mathfrak p}^\times$.
\end{lemma}

\begin{proof}
Map $n$ to the class of $\mathfrak A(n)=\prod_{\mathfrak p}\mathfrak
P_{\mathfrak p}^{n_{\mathfrak p}}$ in $\Cl(K)$ modulo the image of
$\Cl(F)$, a group of order $h_{\rm cap}h_{\rm rel}$. Let $\mathcal N'$ be a
fiber on which the sum of $\mathcal W$ is largest; this sum is positive,
so we can choose $n_0\in\mathcal N'$ with $\mathcal W(n_0)>0$. For
$n\in\mathcal N'$ write $\mathfrak A(n)\mathfrak
A(n_0)^{-1}=x_n\mathfrak a_n\mathcal O_K$ with $x_n\in K^\times$ and
$\mathfrak a_n$ a fractional ideal of $F$, and put
$\beta_n=x_n/\iota x_n$. As in the proof of Lemma~\ref{geo:selection},
$N_{K/F}\beta_n=1$ and $\beta_n$ generates the displayed ideal, which
is supported on the primes $\mathfrak P_{\mathfrak p},\iota\mathfrak
P_{\mathfrak p}$; hence $\beta_n\in\mathcal O_{\mathcal P}$. Two elements of relative
norm one generating the same ideal differ by an element of $\mathcal
O_K^1$, and distinct $n\in\mathcal N'$ give distinct ideals. Finally, if
$\beta\in\beta_n\mathcal O_K^1$, then $\beta_{\mathfrak p}$ has
valuation $n_{\mathfrak p}-n_{0,\mathfrak p}$ and $(\iota\beta)_{\mathfrak
p}=\beta_{\mathfrak p}^{-1}$, which gives the last assertion.
\end{proof}

\subsection{Transfer with Shell Profiles}

This subsection proves Proposition~\ref{fw:transfer}, the form of the
construction used in Section~\ref{sec:certificate}. Its proof follows that of
Proposition~\ref{geo:transfer}, with the window and the lattice extended by
the finite coordinates of Definition~\ref{fw:finite-coordinates}.

\begin{proposition}[Transfer with shell profiles]\label{fw:transfer}
Let $0<\delta<1$ and $C\in\R$. Let $(K_j/F_j)_{j\ge1}$ be quadratic
extensions satisfying (G1)--(G3) with a common $\lambda$, with
$d_j=[F_j:\Q]\to\infty$ and $\log L_{F_j}(1)\le d_jC$. For each $j$ let
$\mathcal P_j$ be a finite set of finite places of $F_j$ that split in
$K_j$, and for $\mathfrak p\in\mathcal P_j$ let $g_{\mathfrak p}$ be a
shell profile (Definition~\ref{fw:shell-profile}) on $F_{\mathfrak p}^2$ with
parameters $(Q_{\mathfrak p},k_{\mathfrak p},w^{(\mathfrak p)})$, where the
parameters range over a fixed finite set independent of $j$. Let
$\mathcal F_{\delta,\mathfrak p}=\mathcal F_\delta(g_{\mathfrak p})$ be the
local functional (Definition~\ref{fw:local-functional}), let $(f_\R,g_\C)$ be
admissible with Fourier constants $(M,\varsigma)$
(Definition~\ref{geo:admissible}), and suppose that
\eqref{fw:fourier-condition} holds for every $j$, with $\mathcal
C_f=\prod_{\mathfrak p\in\mathcal P_j}(\mathcal O_{\mathfrak
p}\times\varpi_{\mathfrak p}^{-k_{\mathfrak p}}\mathcal O_{\mathfrak p})$
the product of the periods of the $g_{\mathfrak p}$, that is, with
$H_f=d_j^{-1}\sum_{\mathfrak p\in\mathcal P_j}k_{\mathfrak p}\log
Q_{\mathfrak p}$. Let $\theta_j$ be the value of $\theta$ for $F_j$, and
put
\begin{equation}\label{fw:margin}
 \begin{split}
 \mathcal M_{{\rm fin},j}={}&
 \frac1{d_j}\sum_{\mathfrak p\in\mathcal P_j}\log\mathcal F_{\delta,\mathfrak p}
 -(1/2-\delta)\ell-C+(1-\delta)\log2+(1-\theta_j)\log\pi\\
 &+(1-2\theta_j)J_\R+\theta_jJ_\C ;
 \end{split}
\end{equation}
we call $\mathcal M_{{\rm fin},j}$ the \emph{margin} of $K_j/F_j$ with
these shell profiles. If $\inf_j\mathcal M_{{\rm fin},j}>0$, then there
are finite sets $U_j\subset\R^2$ with $|U_j|\to\infty$ and
$u(U_j)/|U_j|^{1+\delta}\to\infty$.
\end{proposition}

\begin{remark}\label{fw:margin-comparison}
Compared with \eqref{geo:general-margin}, the margin \eqref{fw:margin}
has no separate term $-\delta H$: by \eqref{fw:shell-kernel}, each
$\mathcal F_{\delta,\mathfrak p}$ contains the factor $Q_{\mathfrak
p}^{-\delta k_{\mathfrak p}}$, and Corollary~\ref{fw:uniform-types}
shows that for uniform local types and unweighted shell profiles
\eqref{fw:margin} reduces to \eqref{geo:general-margin}.
\end{remark}

\begin{proof}[Proof of Proposition~\ref{fw:transfer}]
Fix $\eta>0$ with $4\eta<\inf_j\mathcal M_{{\rm fin},j}$, and consider
one extension $K/F$ of the sequence, with $d$ large. We drop the sequence
index $j$ and write $\mathcal M_{\rm fin}$ for the margin of $K/F$;
below, $j$ indexes the complex places of $F$.

\emph{The profile and the probability measure.} On $\C^d\times\mathbb
A_{\mathcal P}$ let the \emph{total profile} be
\[
 f_{\rm tot}=\prod_vf_\R\prod_{j=1}^cg_\C\prod_{\mathfrak p\in\mathcal
 P}g_{\mathfrak p},
\]
with one factor for each real place $v$ of $F$, each
complex place $j$ of $F$ and each $\mathfrak p\in\mathcal P$, let
$V=-\log f_{\rm tot}$ on the set where $f_{\rm tot}>0$, and put, for the
total mass and the total overlap,
\[
 A_{\rm tot}=\int f_{\rm tot}^p=A_\R^bA_\C^c\prod_{\mathfrak p\in\mathcal P}A(g_{\mathfrak p}),\qquad
 I_{\rm tot}=I_\R^bI_\C^c\prod_{\mathfrak p\in\mathcal P}I(g_{\mathfrak p}).
\]
Besides the points $\mathbf e(u,n)$ defined before Lemma~\ref{fw:selection},
let $\mathbf e(u,n,\omega)$, for $\omega=(\omega_{\mathfrak
p})\in\prod_{\mathfrak p}\mathcal O_{\mathfrak p}^\times$, have
archimedean part $\mathbf e(u)$ and finite part $(\varpi_{\mathfrak
p}^{n_{\mathfrak p}}\omega_{\mathfrak p},\varpi_{\mathfrak
p}^{-n_{\mathfrak p}}\omega_{\mathfrak p}^{-1})_{\mathfrak p}$. Let
$\mathbb P$ be the probability measure on $(\C^d\times\mathbb
A_{\mathcal P})\times\R^c\times\Z^{\mathcal P}\times\prod_{\mathfrak
p}\mathcal O_{\mathfrak p}^\times$ with density
$f_{\rm tot}(x)f_{\rm tot}(x+\mathbf e(u,n,\omega))/I_{\rm tot}$ with respect to $dx\,du\,d\omega$ and
counting measure in $n$, where $d\omega$ is the Haar probability
measure. It is the product of the endpoint laws of the archimedean
blocks (Definition~\ref{geo:profiles}) and of the endpoint laws
(Definition~\ref{fw:endpoint-law}) of the
$g_{\mathfrak p}$. Its endpoints are $x$ and $x+\mathbf e(u,n,\omega)$. Their
energies $V(x)$ and $V(x+\mathbf e(u,n,\omega))$ have the same distribution
and are sums of independent terms, one for each block and each
$\mathfrak p\in\mathcal P$: at the archimedean blocks this is shown in
the proof of Lemma~\ref{geo:concentration}, and at $\mathfrak
p\in\mathcal P$ the map $(z,n,\omega)\mapsto(-z-s(n,\omega),n,\omega)$
preserves the endpoint law of $g_{\mathfrak p}$ and carries its first
endpoint to minus the second, $g_{\mathfrak p}$ is even, and the energy
$-\log g_{\mathfrak p}$ of the first endpoint takes finitely many values.
Let $d\bar m$ be their common mean; this $\bar m$ includes the places of
$\mathcal P$ and is not the $\bar m$ of Lemma~\ref{geo:concentration}.
The $Q_{\mathfrak p}$ take finitely many values, since the parameters of
the shell profiles do, and each is a power of the residue characteristic
of $\mathfrak p$; so these characteristics lie in a fixed finite set,
and $|\mathcal P|=O(d)$. Since the parameters range over a finite set,
the variance of either energy is $O(d)$. Let $\chi$ be the indicator of
the event that both energies lie in $[d(\bar m-\eta),d(\bar m+\eta)]$.
By Chebyshev's inequality, as in the proof of
Lemma~\ref{geo:concentration}, $\mathbb P(\chi=1)\ge1/2$ for large $d$.

\emph{The set $\Omega$.} Let $\Omega$ be the set of points of
$\C^d\times\mathbb A_{\mathcal P}$ where $f_{\rm tot}>0$ and $V\le d(\bar m+\eta)$;
it plays the role of the window of Definition~\fref{geo:window}. It is a
bounded Borel set: the archimedean energies are coercive and bounded
below, and the finite part lies in the compact support of
$\prod_{\mathfrak p}g_{\mathfrak p}$, on which the finite energies are
bounded. It is invariant under rotations of the archimedean
coordinates, and, by Lemma~\ref{fw:shell-invariance}, under
$(x_{\mathfrak p},y_{\mathfrak p})\mapsto(\omega x_{\mathfrak
p},\omega'y_{\mathfrak p})$ for units at each $\mathfrak p\in\mathcal P$
and under translation by $\mathcal C_f$. For $u\in\R^c$ and
$n\in\Z^{\mathcal P}$ let $\Phi(u,n)=|\Omega\cap(\Omega-\mathbf e(u,n))|$, and
let $\mathcal W(n)=\int_{\R^c}\Phi(u,n)\,du$. Only finitely many $n$ have
$\mathcal W(n)\ne0$. By these invariances, the measure of the
intersection of $\Omega$ with its translate by any vector whose
archimedean part has the moduli of $\mathbf e(u)$ and whose finite part is
that of some $\mathbf e(u,n,\omega)$ equals $\Phi(u,n)$. Where $\chi=1$, both
endpoints lie in $\Omega$ and $f_{\rm tot}(x)f_{\rm tot}(x+\mathbf e(u,n,\omega))\le
e^{-2d(\bar m-\eta)}$. Integrating over $x$, $u$ and $\omega$ and
summing over $n$, we obtain
\[
 \begin{split}
 \sum_n\mathcal W(n)
 &=\sum_n\int\mathbf 1_\Omega(x)\,\mathbf 1_\Omega(x+\mathbf e(u,n,\omega))\,
 dx\,du\,d\omega\\
 &\ge e^{2d(\bar m-\eta)}\sum_n\int\chi\,f_{\rm tot}(x)f_{\rm tot}(x+\mathbf e(u,n,\omega))\,
 dx\,du\,d\omega\\
 &=e^{2d(\bar m-\eta)}I_{\rm tot}\,\mathbb P(\chi=1)
 \ge\tfrac12I_{\rm tot}e^{2d(\bar m-\eta)} .
 \end{split}
\]
Moreover $\mathbf 1_\Omega\le e^{pd(\bar m+\eta)}f_{\rm tot}^p$, so
$|\Omega|\le A_{\rm tot}e^{pd(\bar m+\eta)}$.

\emph{Class selection and rescaling.} Apply Lemma~\ref{fw:selection} to
the finitely many $n$ with $\mathcal W(n)>0$, obtaining $n_0$,
$\mathcal N'$ and $\beta_n$, and let $\varpi_0$ be as in that lemma.
Put $\Omega'=\varpi_0^{-1}\Omega$, where $\varpi_0$ acts on the finite
part. Then $|\Omega'|=|\Omega|$, and $\Omega'$ is invariant under the
rectangular group $\mathcal C'_f=\varpi_0^{-1}\mathcal C_f$, of measure
$|\mathcal C_f|$. For $n\in\mathcal N'$, $\beta\in\beta_n\mathcal O_K^1$
and $h\in\R^c$, the element $\beta^{(h)}=(W_h\beta_\infty,\beta_f)\in
\mathcal O_{\mathcal P,h}$ satisfies
\[
 |\Omega'\cap(\Omega'-\beta^{(h)})|=|\Omega\cap(\Omega-(W_h\beta_\infty,
 \varpi_0\beta_f))|=\Phi(l(\beta)-h,n),
\]
by the last assertion of Lemma~\ref{fw:selection} and the invariances of
$\Omega$, as in Lemma~\ref{geo:phi-omega}.

\emph{Points and edges.} Let $\mathfrak a_{\mathcal C'_f}$ be the ideal attached to
$\mathcal C'_f$ in Lemma~\ref{fw:lattice}. Fix a $\Z$-basis of
$(\mathfrak a_{\mathcal C'_f})_\infty$, let $P_{\mathcal C'_f}$ be the fundamental
parallelepiped that it spans, and for $y\in[0,1)^{2d}$ let
$y_{\mathcal C'_f}\in P_{\mathcal C'_f}$ be the point with coordinate
vector $y$ in this basis. For $h\in\R^c$, $y\in[0,1)^{2d}$ and
$\zeta\in\mathcal C'_f$ put $z_0=(W_hy_{\mathcal C'_f},\zeta)$; these
points form the fundamental domain $W_hP_{\mathcal C'_f}\times\mathcal
C'_f$ of $\mathcal O_{\mathcal P,h}$ of Lemma~\ref{fw:lattice}(d). Let
$\mathcal X(h,z_0)=(\mathcal O_{\mathcal P,h}+z_0)\cap\Omega'$ and let
$N(h,z_0)=|\mathcal X(h,z_0)|$ be the number of its points (the analogue
of the count $n(h,y)$ of Definition~\fref{geo:counts}; we write $N$ because $n$
indexes $\Z^{\mathcal P}$ in this proof), and let
$\mathcal E(h,z_0)$ count the pairs $(x,\beta)$ with
$x\in\mathcal X(h,z_0)$, $\beta\in\bigcup_{n\in\mathcal N'}\beta_n\mathcal
O_K^1$ and $x+\beta^{(h)}\in\Omega'$. Averages over $(h,z_0)$ are taken
with respect to the probability measure
$dh\,dy\,d\zeta/(\operatorname{Reg}^1|\mathcal C'_f|)$ on
$\Pi^1\times[0,1)^{2d}\times\mathcal C'_f$. Since
$\mathbf 1_{\Omega'}\le e^{pd(\bar m+\eta)}(f_{\rm tot}\circ\varpi_0)^p$ and
$(f_{\rm tot}\circ\varpi_0)^p$ is the product of an archimedean factor and a
function invariant under $\mathcal C'_f$, Lemmas~\ref{fw:counting}
and~\ref{fw:rescaling} give
\[
 N(h,z_0)\le N_{\max}=\frac{2A_{\rm tot}e^{pd(\bar m+\eta)}}{\covol_0}
\]
for all $(h,z_0)$. By tiling, for fixed $h$ the average of $N(h,z_0)$
over $(y,\zeta)$ is $|\Omega'|/\covol_0$, and the average of
$\mathcal E(h,z_0)$ is
$\covol_0^{-1}\sum_{n\in\mathcal N'}\sum_{\beta\in\beta_n\mathcal O_K^1}
\Phi(l(\beta)-h,n)$. Averaging over $h\in\Pi^1$ as in
Lemma~\ref{geo:unfold-lemma}, the average of $\mathcal E$ becomes
\[
 \frac{w_K}{\operatorname{Reg}^1\covol_0}\sum_{n\in\mathcal N'}\mathcal W(n)
 \ge\frac{w_K}{h_{\rm cap}h_{\rm rel}\operatorname{Reg}^1\covol_0}
 \sum_n\mathcal W(n).
\]
Measurability and integrability hold as in
Lemma~\ref{geo:joint-average-lemma}, since $\mathcal E\le N^2\le
N_{\max}^2$. The argument of that lemma yields a pair $(h,z_0)$ with
$N=N(h,z_0)>0$ and
\[
 \frac{\mathcal E(h,z_0)}N\ge\frac{w_K}{h_{\rm cap}h_{\rm rel}
 \operatorname{Reg}^1}\,\frac{\sum_n\mathcal W(n)}{|\Omega|}
 \ge\frac{w_K}{h_{\rm cap}h_{\rm rel}\operatorname{Reg}^1}\,
 \frac{I_{\rm tot}e^{2d(\bar m-\eta)}}{2A_{\rm tot}e^{pd(\bar m+\eta)}} .
\]
With $N\le N_{\max}$ and $p(1+\delta)=2$,
\[
 \frac{\mathcal E(h,z_0)}{N^{1+\delta}}\ge
 \frac{w_K\covol_0^\delta}{2^{1+\delta}h_{\rm cap}h_{\rm rel}
 \operatorname{Reg}^1}\,\frac{I_{\rm tot}}{A_{\rm tot}^{1+\delta}}\,e^{-4\eta d}.
\]
By \eqref{geo:mass-density} and $\log L_F(1)\le dC$, the first factor
is at least
$2^{-2-\delta}2^{d(1-\delta)}\pi^{(1-\theta)d}\lambda^{-(1/2-\delta)d}e^{-dC}$,
and $I_{\rm tot}/A_{\rm tot}^{1+\delta}=\exp\bigl(bJ_\R+cJ_\C+\sum_{\mathfrak p}\log\mathcal
F_{\delta,\mathfrak p}\bigr)$. Hence
$\mathcal E/N^{1+\delta}\ge2^{-2-\delta}e^{d(\mathcal M_{\rm
fin}-4\eta)}$.

\emph{The planar set.} Let $v_0$ be a real place of $F$ and let $U$ be
the set of $v_0$-coordinates of the points of $\mathcal X(h,z_0)$.
Distinct points differ by $(W_h\gamma_\infty,\gamma_f)$ with
$0\ne\gamma\in\mathcal O_{\mathcal P}$, whose $v_0$-coordinate is nonzero, so $|U|=N$.
Each counted pair $(x,\beta)$ gives two points at distance
$|\phi_{v_0}(\beta)|=1$, and distinct pairs give distinct ordered
pairs of points. As in the proof of Proposition~\ref{geo:transfer}
(Lemma~\ref{geo:planar}), $u(U)\ge\mathcal E/2$, the ratio
$u(U)/|U|^{1+\delta}$ tends to infinity, and the bound $u(U)\le|U|^2/2$
forces $|U|\to\infty$.
\end{proof}

\subsection{Uniform Local Types}

This subsection shows that Proposition~\ref{fw:transfer} contains
Proposition~\ref{geo:transfer} when the places carry uniform local types. The
proof of Theorem~\ref{thm:main} does not use Corollary~\ref{fw:uniform-types},
since the types at $41$ are not uniform
(Definition~\fref{cert:effective-types}); it applies
Proposition~\ref{fw:transfer} directly.

\begin{corollary}[Uniform local types]\label{fw:uniform-types}
Let $0<\delta<1$ and $C\in\R$. Let $(K_j/F_j)_{j\ge1}$ be quadratic
extensions satisfying (G1)--(G3) with a common $\lambda$, with
$d_j=[F_j:\Q]\to\infty$ and $\log L_{F_j}(1)\le d_jC$, and with the same
uniform local types $(e_r,f_r)_{r\in\mathcal R}$
(Definition~\ref{geo:uniform-types}). For each $r\in\mathcal R$ fix an
integer $k_r\ge0$ and weights $w^{(r)}$ as in
Definition~\ref{fw:shell-profile}, let $\mathcal F_{\delta,r}$ be the
value of the local functional $\mathcal F_\delta$ at the shell profile
with parameters $(r^{f_r},k_r,w^{(r)})$, and let $H$ be as in
\eqref{geo:HJ}. Let $(f_\R,g_\C)$ be admissible (Definition~\ref{geo:admissible}) with Fourier constants
$(M,\varsigma)$ satisfying \eqref{geo:poisson-condition} with
$\mu=2e^{H/2-\ell}$, and put
\[
 \begin{split}
 \mathcal M_{\mathcal R}(\theta)={}&
 \sum_{r\in\mathcal R}\frac{\log\mathcal F_{\delta,r}}{e_rf_r}
 -(1/2-\delta)\ell-C+(1-\delta)\log2+(1-\theta)\log\pi\\
 &+(1-2\theta)J_\R+\theta J_\C .
 \end{split}
\]
If $\inf_j\mathcal M_{\mathcal R}(\theta_j)>0$, where $\theta_j$ is the
value of $\theta$ for $F_j$, then there are finite sets
$U_j\subset\R^2$ with $|U_j|\to\infty$ and
$u(U_j)/|U_j|^{1+\delta}\to\infty$. For unweighted shell profiles,
$\sum_r\log\mathcal F_{\delta,r}/(e_rf_r)=J-\delta H$, so that
$\mathcal M_{\mathcal R}$ is the margin $\mathcal M$ of
\eqref{geo:general-margin}.
\end{corollary}

\begin{proof}
Apply Proposition~\ref{fw:transfer} with $\mathcal P_j$ the set of all
places of $F_j$ above $\mathcal R$, which split in $K_j$, and with
$g_{\mathfrak p}$ the shell profile with parameters
$(r^{f_r},k_r,w^{(r)})$ for $\mathfrak p$ above $r$; these parameters
take $|\mathcal R|$ values, and $\mathcal F_{\delta,\mathfrak
p}=\mathcal F_{\delta,r}$. By Lemma~\ref{geo:prime-count} there are $d/(e_rf_r)$
places above $r$, each with $Q_{\mathfrak p}=r^{f_r}$, so
\begin{equation}\label{fw:uniform}
 \frac1d\sum_{\mathfrak p\in\mathcal P}\log\mathcal F_{\delta,\mathfrak p}
 =\sum_{r\in\mathcal R}\frac{\log\mathcal F_{\delta,r}}{e_rf_r},
 \qquad
 H_f=\sum_{r\in\mathcal R}\frac{k_r\log r}{e_r}=H .
\end{equation}
In particular \eqref{fw:fourier-condition} coincides with
\eqref{geo:poisson-condition}, whatever the shell weights, and the
margin \eqref{fw:margin} of $K_j/F_j$ equals $\mathcal M_{\mathcal
R}(\theta_j)$. For unweighted shell profiles, Lemma~\ref{fw:shell-lemma}
gives $\log\mathcal F_{\delta,r}=\log(k_r+1)-\delta k_rf_r\log r$, and the
sum in \eqref{fw:uniform} is $J-\delta H$.
\end{proof}

Thus Proposition~\ref{fw:transfer} contains
Proposition~\ref{geo:transfer}, and the shell weights replace the term
$J-\delta H$ of \eqref{geo:general-margin} by
$\sum_r\log\mathcal F_{\delta,r}/(e_rf_r)$ without changing the Fourier
condition.
